\documentclass[10pt,twocolumn]{article}
\usepackage[a4paper,margin=18mm,columnsep=7mm]{geometry}
\usepackage{amsmath,amssymb,mathtools,bm,microtype}
\usepackage[hidelinks]{hyperref}

\newcommand{\F}{\mathbb F_2}
\newcommand{\MWPM}{\operatorname{MWPM}}

\title{Circuit-level extension of the color-code swim distance}
\author{}
\date{}

\begin{document}
\maketitle

\section{Circuit-level extension}
\label{sec:circuit_level_main}

The perfect-measurement construction associates, for each color \(c\), a
modified stage-2 decoding graph whose two inequivalent boundaries encode the
two logical classes within a fixed stage-1 fiber.  The swim distance is then
the minimum residual distance between these boundaries after contracting the
clusters produced by the decoder.  Here we extend this construction to a
fully terminated circuit-level memory experiment.  The main difference is
that, at circuit level, a logical class is most naturally defined directly
from the detector error model (DEM), rather than from a particular spacetime
drawing.

\subsection{Fixed-fiber logical classes from the DEM}

We consider one Pauli sector and one color \(c\).  After the
Lee--Li--Bartlett decomposition, the stage-2 \(c\)-only DEM is graphlike and
contains physical \(c\)-colored detector rows, virtual detector rows inherited
from stage 1, and logical-observable labels \cite{Lee}.  Let
\[
 D_c^{\rm circ}:\F^{E_c^{\rm circ}}\rightarrow\F^{V_c^{\rm real}},
 \qquad
 L_c^{\rm circ}:\F^{E_c^{\rm circ}}\rightarrow\F
\]
denote respectively the detector-incidence and logical-observable maps of the
retained stage-2 model.  A binary vector
\(x\in\F^{E_c^{\rm circ}}\) specifies a set of retained stage-2 error
mechanisms, so that
\[
 s=D_c^{\rm circ}x,\qquad
 \ell=L_c^{\rm circ}x
\]
are its stage-2 syndrome and logical frame change.

Fixing the stage-1 output fixes the virtual detector values and therefore a
stage-2 fiber
\[
 \mathcal F_s=\{x:D_c^{\rm circ}x=s\}.
\]
Two candidates \(x,x'\in\mathcal F_s\) differ by
\(z=x+x'\in\ker D_c^{\rm circ}\).  Their logical frames agree when
\(L_c^{\rm circ}z=0\), and are opposite when
\(L_c^{\rm circ}z=1\).  Thus the circuit-level analogue of the
perfect-measurement logical quotient is
\[
 \frac{\ker D_c^{\rm circ}}
 {\ker D_c^{\rm circ}\cap\ker L_c^{\rm circ}}
 \cong
 \operatorname{im}
 \!\left(
 L_c^{\rm circ}|_{\ker D_c^{\rm circ}}
 \right).
 \label{eq:circ_quotient}
\]
For the one-observable memory experiment this quotient contains two classes
whenever \(L_c^{\rm circ}\) is nonzero on
\(\ker D_c^{\rm circ}\).  Equivalently,
\[
 L_c^{\rm circ}\notin\operatorname{rowspan}D_c^{\rm circ}.
 \label{eq:rank_condition}
\]
This condition depends only on the fixed decoding graph and can therefore be
checked once when the circuit-level decoder is constructed.

This DEM-based formulation avoids identifying temporal boundaries, matching
boundaries, and logical boundaries by coordinates alone.  Detector time,
color, and position remain useful metadata, but the logical action of a
mechanism is determined by \(L_c^{\rm circ}\).

\subsection{Cluster contraction and the circuit-level swim distance}

Let \(w_e\geq0\) be the edge weight of a retained stage-2 mechanism and let
the MWPM decoder provide a set of growth radii.  As in the
perfect-measurement construction, these radii define metric balls on the
original weighted stage-2 graph.  If \(h_u\) and \(h_v\) are the maximum
covered lengths arriving at the two endpoints of an edge
\(e=uv\), the uncovered part of the edge is
\[
 \bar w_e=\max\{0,w_e-h_u-h_v\}.
 \label{eq:residual_weight}
\]
The same interval construction applies to spatial, timelike, and diagonal
circuit-level edges.  Importantly, the coverage is computed on the original
stage-2 graph and is only then transferred to the logical analysis graph.
No new decoder growth is performed after introducing a logical cut or cover.

We define the circuit-level swim distance by
\[
 \boxed{
 \phi_c^{\rm circ}
 =
 \min_{\substack{
 D_c^{\rm circ}z=0\\
 L_c^{\rm circ}z=1
 }}
 \sum_{e}\bar w_e z_e
 }.
 \label{eq:circuit_swim}
\]
Hence \(\phi_c^{\rm circ}\) is the minimum residual cost required to move
from the decoder's current fixed-fiber logical class to the opposite class
without changing the stage-2 syndrome.  This definition is independent of a
particular geometric embedding.

\subsection{General computation by a logical double cover}

To compute Eq.~\eqref{eq:circuit_swim}, complete every one-detector
mechanism at a single artificial matching-boundary vertex \(b_\ast\);
two-detector mechanisms remain ordinary edges, and a zero-detector mechanism
is represented as a labelled loop at \(b_\ast\).  Denote the resulting
multigraph by \(G_c^\bullet\).  Each edge carries the logical label
\[
 \lambda_e=L_c^{\rm circ}[e]\in\F.
\]

We then construct a binary logical cover with two copies of every vertex,
\((v,0)\) and \((v,1)\).  For every edge \(e=uv\), add
\[
 (u,a)\longleftrightarrow(v,a+\lambda_e),
 \qquad a\in\F,
 \label{eq:cover_edge}
\]
and assign both lifted edges the same residual cost \(\bar w_e\).  An
even-labelled edge stays within one sheet, whereas an odd-labelled edge
switches sheets.  A detector-invisible chain with odd logical parity therefore
lifts to a path between opposite copies of the same base vertex.  Consequently,
\[
 \phi_c^{\rm circ}
 =
 \min_{v\in V(G_c^\bullet)}
 \operatorname{dist}_{\widehat G_c}\bigl((v,0),(v,1)\bigr).
 \label{eq:cover_distance}
\]
The minimizing path can be projected back to the original stage-2 mechanisms
and reduced modulo two to obtain a witness \(z\) satisfying
\[
 D_c^{\rm circ}z=0,\qquad
 L_c^{\rm circ}z=1.
\]
This construction is completely graph-based and does not require an explicit
spacetime surface.

\subsection{Two-boundary form for closed memory experiments}

For the standard fully terminated memory experiment, Eq.~\eqref{eq:cover_distance}
often admits a simpler two-boundary realization.  Consider the subgraph
formed by the real vertices and the edges joining two real vertices.  If
every cycle of this internal graph has even logical label, there exists a
binary potential \(g\) such that
\[
 \lambda_{uv}=g(u)+g(v)
 \label{eq:balanced}
\]
for every internal edge.  This condition is checked efficiently by a single
spanning-forest traversal and is therefore an offline preprocessing step,
not a per-shot computation.

Under Eq.~\eqref{eq:balanced}, we gauge-transform
\[
 \lambda'_e
 =
 \lambda_e + g^{\mathsf T}D_c^{\rm circ}[:,e],
\]
which leaves \(L_c^{\rm circ}z\) unchanged for every
\(z\in\ker D_c^{\rm circ}\).  All internal edges can then be made
logical-even.  Splitting the single artificial matching boundary into two
vertices \(b^0\) and \(b^1\), and attaching each boundary edge according to
its transformed logical label, gives a resolved graph
\(G_c^{\rm cut}\) for which
\[
 \boxed{
 \phi_c^{\rm circ}
 =
 \operatorname{dist}_{\bar G_c^{\rm cut}}(b^0,b^1)
 }.
 \label{eq:two_boundary}
\]
Thus, when the balance condition holds, the circuit-level computation has
exactly the same operational form as the perfect-measurement construction:
contract the decoder clusters and measure the shortest residual distance
between two inequivalent boundaries.

The condition is not implied merely by the phrase ``closed temporal
boundary''.  A closed memory circuit can still contain one-detector
mechanisms near initialization or readout.  What matters is whether the
logical parity of every internal cycle is trivial.  If this condition fails,
Eq.~\eqref{eq:cover_distance} remains valid and provides the general fallback.

\subsection{Spacetime interpretation}

The preceding construction also clarifies the relation to the usual
spacetime picture.  A repeated triangular color-code memory can be viewed
as a triangular-prism-like spacetime protocol, with logical information
propagated by correlation surfaces.  In this interpretation,
\(L_c^{\rm circ}\) is the algebraic intersection functional associated with
the relevant logical correlation.  A valid two-boundary cut can be viewed
as opening the decoding graph along a representative of this logical
cochain.

This geometric picture is useful for intuition, but the DEM formulation is
more general: hook errors, diagonal edges, virtual detectors, and compressed
mechanisms need not form a literal Cartesian product of the static decoding
graph with time.  The logical-cover construction remains valid without such
an embedding.

\subsection{Relation to the perfect-measurement result}

The circuit-level definition reduces exactly to the code-capacity result in
the controlled perfect-measurement limit.  If the only noisy mechanisms are
independent data-qubit errors before an ideal syndrome-extraction round, then
the circuit-level detector map reduces to the perfect-measurement stage-2
map and, on detector-invisible differences,
\[
 D_c^{\rm circ}=D_c,\qquad
 L_c^{\rm circ}z=\beta_c(z).
\]
With the same cluster coverage,
\[
 \phi_c^{\rm circ}
 =
 \phi_c,
\]
where \(\phi_c\) is the previously defined distance between the two
resolved perfect-measurement boundaries.  Hence
Eq.~\eqref{eq:circuit_swim} is a genuine extension rather than a separate
confidence definition.

Finally, if the clusters are generated from an exact MWPM optimum together
with the required optimal nonnegative dual certificate, the perfect-measurement
representative-gap argument also extends:
\[
 W_{c,\mathrm{opp}}^{\rm circ}
 -
 W_{c,\mathrm{base}}^{\rm circ}
 \geq
 \phi_c^{\rm circ}.
 \label{eq:circ_gap_bound}
\]
This is a fixed-fiber minimum-representative bound, not a full logical-class
LLR or a theorem about the complete three-color decoder.  In the current
implementation the exported Sparse-Blossom growth data are used as a
practical soft-output convention but are not yet certified as the exact
dual required by Eq.~\eqref{eq:circ_gap_bound}.

The present construction assumes a fully terminated memory experiment.
Removing future detector rows changes \(\ker D_c^{\rm circ}\) and can change
its logical image; therefore open temporal boundaries, including
sliding-window decoding, require a separate treatment.

\begin{thebibliography}{9}

\bibitem{Lee}
S.-H. Lee, A. Li, and S. D. Bartlett,
``Color code decoder with improved scaling for correcting circuit-level
noise,'' \emph{Quantum} \textbf{9}, 1609 (2025).

\bibitem{Meister}
N. Meister, C. A. Pattison, and J. Preskill,
``Efficient soft-output decoders for the surface code,''
arXiv:2405.07433.

\end{thebibliography}

\end{document}
